% Comparison of the two stem models, written for the supervisors (standalone explanation).
% Follows docs/stem_model_tradeoff.tex (option (d) there). Source of the numbers: docs/notes/2026-10-05.md.
%
% This file compiles on its own. To use it inside another document, copy everything between
% "BEGIN CONTENT" and "END CONTENT". The content needs the packages amsmath and booktabs.
\documentclass[11pt,a4paper]{article}
\usepackage[T1]{fontenc}
\usepackage[utf8]{inputenc}
\usepackage[margin=2.5cm]{geometry}
\usepackage{amsmath}
\usepackage{booktabs}

\begin{document}

% ======================================= BEGIN CONTENT =======================================
\section{Stem model comparison: Newton cable and rigid-segment chain}

\subsection{Summary}

The previous section showed that the Newton cable cannot combine realistic stretch and shear stiffness,
correct bending and an affordable simulation. I have therefore built the alternative suggested there, option
(d): a chain of rigid segments whose joints can only bend and twist, simulated as an articulation in PhysX,
the physics engine that also simulates the robot arm. Both models are kept and can be selected by name.

On the same placeholder stem and the same tests, the chain reproduces the thin-rod reference within 1 to
2\,\% (sag, push, first natural frequency), cannot stretch or shear, and reports the force in every joint
exactly. It needed one modelling addition, a small extra inertia on every joint (``armature''), which makes the
second and third natural frequencies too low. It is about 1.8 times slower than the cable at its current
setting. Contact with the robot's tool has not been tested yet.

\subsection{The chain model}

The stem is again 20 rigid segments of 2\,cm (capsules of 8\,mm diameter). The lowest segment is clamped to
the ground by a fixed joint. Every other joint connects two neighbouring segments at their common end and can
only rotate, about three axes: two bending axes and the stem axis (twist). Each rotation has a spring and a
damper with the same values as in the cable,
\begin{equation*}
  k_\text{bend} = \frac{EI}{l}, \qquad k_\text{twist} = \frac{GJ}{l}, \qquad
  c = \tau\, k \quad (\tau = 3.2\,\text{ms}),
\end{equation*}
where $l$ is the segment length. There are no stretch or shear springs: the joints do not allow these motions
at all, which is the thin-rod assumption. The model is described by its joint angles (reduced coordinates),
like a robot arm, and runs in the same solver as the arm, so robot and stem need no coupling of two solvers.

\subsection{First finding: the joint springs are also solved iteratively}

In the first version the chain was far too soft: under its own weight the horizontal stem sagged 16 times more
than expected, and the upright stem fell over after a kick. The reason is similar to the cable's problem. PhysX
solves the joint springs iteratively, a fixed number of passes per time step. The segments are very light
(1\,g) and the springs comparatively stiff, so the chain has natural vibrations up to 3.6\,kHz, in which single
segments wobble against their neighbours. A few passes per time step cannot resolve them, and the springs do
not reach their equilibrium. More solver work helps only slowly (Table~\ref{tab:chain-iterations}); what counts
is the total number of passes per simulated second, not the time step alone.

\begin{table}[h]
  \centering
  \caption{Chain without armature: sag of the horizontal stem under its own weight (expected 14.17\,mm).}
  \label{tab:chain-iterations}
  \begin{tabular}{lll}
    \toprule
    Time step & Passes per step & Sag \\
    \midrule
    2\,ms   & 8   & 228\,mm, stem falls over \\
    2\,ms   & 32  & 225\,mm \\
    2\,ms   & 128 & 62\,mm \\
    0.5\,ms & 32  & 61\,mm \\
    \bottomrule
  \end{tabular}
\end{table}

\subsection{The fix: joint armature}

PhysX offers ``armature'', a constant rotational inertia added to every joint coordinate. The name comes from
motors, where the rotor inertia appears at the joint through the gearbox. A natural frequency is
$\omega = \sqrt{k / I_\text{eff}}$. In the fast modes a single segment moves against its neighbours; its own
inertia is tiny (about $10^{-7}$\,kg\,m$^2$), so an armature of $10^{-4}$\,kg\,m$^2$ increases the effective
inertia about a thousandfold and lowers these frequencies from 3.6\,kHz to 36\,Hz. In the first mode the whole
stem swings; its inertia is already large (about $9\cdot10^{-4}$\,kg\,m$^2$ about the base), so the same
armature changes it by only 1\,\%. The ratio between the fastest and the slowest mode, which governs how hard the
problem is for an iterative solver, drops from about 700 to about 7.

The static shape (sag, deflection under a steady push) does not depend on inertia at all and is unaffected. The
price is that fast motions become too slow. I computed the natural frequencies of the chain exactly, from a
linear small-deflection model of the same segments and springs (Table~\ref{tab:chain-armature}).

\begin{table}[h]
  \centering
  \caption{Natural frequencies of the 20-segment chain (exact linear model, no gravity) and simulated sag.
           For comparison, a continuous clamped beam has 5.20, 32.6 and 91.2\,Hz.}
  \label{tab:chain-armature}
  \begin{tabular}{lllll}
    \toprule
    Armature [kg\,m$^2$] & Mode 1 & Mode 2 & Mode 3 & Simulated sag (expected 14.17\,mm) \\
    \midrule
    0                  & 5.21\,Hz          & 32.7\,Hz          & 91.5\,Hz          & 228\,mm \\
    $10^{-5}$          & 5.20\,Hz          & 31.4\,Hz          & 71.1\,Hz          & 16.2\,mm \\
    $3\cdot10^{-5}$    & 5.19\,Hz          & 29.2\,Hz (-11\,\%) & 53.1\,Hz (-42\,\%) & 14.8\,mm \\
    $10^{-4}$ (in use) & 5.15\,Hz (-1\,\%) & 24.1\,Hz (-26\,\%) & 33.2\,Hz (-64\,\%) & 14.3\,mm \\
    \bottomrule
  \end{tabular}
\end{table}

The time step matters as well (Table~\ref{tab:chain-dt}). With 2\,ms all tests are within 2\,\%; with 10\,ms
the push is 12\,\% too large.

\begin{table}[h]
  \centering
  \caption{Chain with armature $10^{-4}$\,kg\,m$^2$ and 8 passes per step. Expected: push 9.44\,mm,
           sag 14.17\,mm, first frequency 5.06\,Hz (exact, with gravity and armature).}
  \label{tab:chain-dt}
  \begin{tabular}{lllll}
    \toprule
    Time step & Push & Sag & First frequency & Damping ratio \\
    \midrule
    2\,ms (in use) & 9.60\,mm (+1.7\,\%) & 14.30\,mm (+0.9\,\%) & 5.00\,Hz & 0.056 \\
    5\,ms          & 9.90\,mm (+4.9\,\%) & 14.26\,mm (+0.6\,\%) & 4.92\,Hz & 0.063 \\
    10\,ms         & 10.55\,mm (+12\,\%) & 13.53\,mm (-4.5\,\%) & 4.76\,Hz & 0.072 \\
    \bottomrule
  \end{tabular}
\end{table}

\subsection{Comparison}

Both models use the same placeholder stem (length 0.4\,m, diameter 8\,mm, density 1000\,kg/m$^3$, bend modulus
$5\cdot10^{8}$\,Pa, 20 segments), clamped at the base, without robot or contact. The tests are those of the
previous section: sag of the horizontal stem, push of 0.05\,N on the last segment of the upright stem, and a
kick (everything above the lowest joint is set rotating so that the top moves sideways at 1\,m/s). The reference
is the thin-rod chain: bending only, computed exactly for the same 20 segments, including gravity where it
matters.

\begin{table}[h]
  \centering
  \caption{Cable (Newton, stretch 0.01 and shear 0.001 times the bend modulus, cost 80, 10\,ms step) and chain
           (PhysX, armature $10^{-4}$\,kg\,m$^2$, 2\,ms step, 8 passes).}
  \label{tab:model-comparison}
  \begin{tabular}{llll}
    \toprule
    Test & Thin-rod reference & Cable & Chain \\
    \midrule
    Sag under own weight      & 14.17\,mm   & 15.78\,mm (+11\,\%)    & 14.30\,mm (+0.9\,\%) \\
    Push 0.05\,N, upright     & 9.44\,mm    & 10.2\,mm (+8\,\%)      & 9.60\,mm (+1.7\,\%) \\
    First natural frequency   & 5.11\,Hz    & 4.67\,Hz (-9\,\%)      & 5.00\,Hz (-2\,\%) \\
    Damping ratio             & about 0.05  & 0.063                  & 0.056 \\
    Compression under own weight & practically none & 2.3 times the value of its own springs & none \\
    Force in the joints       & --          & not available          & exact \\
    Higher natural frequencies & 32.7, 91.5\,Hz & not measured        & 24.1, 33.2\,Hz (model) \\
    Same result in every parallel environment & -- & only if all are placed at the origin & yes \\
    Computing time, 1024 stems, per 10\,ms & -- & 8.6\,ms          & about 15.8\,ms \\
    Solvers needed with the robot & --      & two, coupled           & one \\
    Contact force readable    & --          & not with Isaac Lab's sensor & expected, not tested \\
    \bottomrule
  \end{tabular}
\end{table}

The cable reproduces its own softened model well (within 1\,\% of a hand calculation that includes its soft
shear spring), but that model is about 10\,\% softer than the thin rod. The chain reproduces the thin rod itself.
The computing time of the chain is a rough measurement that includes per-step overhead of the test script.

\subsection{Damage measures in the chain}

Because the joints of the chain are ideal constraints, the solver knows the force that each joint transmits. This
gives the axial force and the twisting moment in every joint directly, which is what the tearing and torsion
limits need (Table~\ref{tab:chain-forces}). In the cable the axial force can only be estimated from the
stretch of the softened springs, which has not converged at the current solver cost.

\begin{table}[h]
  \centering
  \caption{Forces read from the joints of the upright chain at rest (positive = compression). Expected axial
           force: the weight of the segments above the joint, minus the pull.}
  \label{tab:chain-forces}
  \begin{tabular}{llll}
    \toprule
    Load & Joint 1 (bottom) & Joint 10 & Joint 19 (top) \\
    \midrule
    Own weight                  & 0.1874\,N  & 0.0986\,N  & 0.0099\,N \\
    Expected                    & 0.1874\,N  & 0.0986\,N  & 0.0099\,N \\
    Own weight + 1\,N pull up   & -0.8126\,N & -0.9014\,N & -0.9901\,N \\
    \midrule
    Twisting moment, 0.01\,N\,m on the top segment & 0.0100\,N\,m & & \\
    \bottomrule
  \end{tabular}
\end{table}

One caveat concerns local damage. In the chain the tool always touches a rigid segment of 2\,cm, so a local
twisting or pressing load is spread over that segment and its two joints, while a real stem can be damaged
locally. A policy might learn to twist or press a segment hard without being penalized by a curvature limit.
The damage limits should therefore also bound the twisting moment and the contact force, not only the
curvature.

\subsection{Limits of these results}

\begin{itemize}
  \item One stem geometry with placeholder values, the stem alone: no robot, no contact with the tool.
  \item The armature value was chosen for this geometry. Lighter or shorter segments (randomized diameter or
        length) need it checked again.
  \item Only slow motions are reproduced accurately; the second and third natural frequencies of the chain are
        26\,\% and 64\,\% too low.
\end{itemize}

\subsection{Questions for discussion}

\begin{enumerate}
  \item Is joint armature an accepted way to stabilize such a chain, given that it distorts the higher natural
        frequencies? For slow pushing the static shape and the first mode should matter most.
  \item Would a solver that treats the joint dynamics directly instead of iteratively be preferable? The
        MuJoCo-based solver available in Isaac Lab solves the joint-space equations with a direct method and
        may not need armature (untested).
  \item With the chain, robot and stem run in one solver, and contact forces should be readable. Is this the
        model to continue with for the first training stage, keeping the cable for comparison?
  \item Which damage limits besides curvature are needed: twisting moment, axial force (tearing), contact force
        (crushing)? Are realistic values for plant stems known?
\end{enumerate}
% ======================================== END CONTENT ========================================

\end{document}
